\chapter{Rekonstruktion aus Mengenprodukten}

\section{Die große Potenzstruktur einer Gruppe}

\begin{definition}[Nichtleere Teilmengen mit Mengenprodukt]
Für eine Gruppe \(G\) sei
\[
 \mathcal P_*(G)=\{A\subseteq G:A\ne\varnothing\},\qquad
 AB=\{ab:a\in A,\ b\in B\}.
\]
Ein \emph{Isomorphismus der Mengenproduktstrukturen}
ist eine Bijektion \(F:\mathcal P_*(G)\to\mathcal P_*(H)\) mit
\(F(AB)=F(A)F(B)\).
Die Struktur \((\mathcal P_*(G),\cdot)\) heißt auch
\emph{große Potenzhalbgruppe} von \(G\).
Hier werden ausschließlich ihre gruppenspezifischen
Rekonstruktionseigenschaften betrachtet.
\end{definition}

\begin{theorem}[Einermengen sind genau die invertierbaren Teilmengen]
\label{b40guide:thm:singletons}
Die Einheit des Mengenprodukts auf \(\mathcal P_*(G)\) ist \(\{1\}\).
Für \(A\in\mathcal P_*(G)\) sind äquivalent:
\begin{enumerate}
\item Es gibt \(B\in\mathcal P_*(G)\) mit \(AB=BA=\{1\}\).
\item \(A=\{g\}\) für ein \(g\in G\).
\end{enumerate}
Im zweiten Fall ist die inverse Teilmenge \(\{g^{-1}\}\).
Mit dem Mengenprodukt bilden die Einermengen eine Gruppe, und
\(g\mapsto\{g\}\) ist ein Isomorphismus von \(G\) auf diese Gruppe.
\cite{projectreconstruction,liutringali}
\end{theorem}

\begin{theorem}[Erkennung des vollen Gruppenträgers]
Für jede nichtleere Teilmenge \(A\subseteq G\) gilt
\[
 AG=GA=G.
\]
Die Menge \(G\) ist das einzige \(Z\in\mathcal P_*(G)\) mit
\(AZ=ZA=Z\) für alle \(A\in\mathcal P_*(G)\).
Daher erfüllt jeder Isomorphismus
\(F:\mathcal P_*(G)\to\mathcal P_*(H)\) die Gleichungen~\cite{projectreconstruction}
\[
 F(\{1_G\})=\{1_H\},\qquad F(G)=H.
\]
\end{theorem}

\section{Rückgewinnung der Gruppe}

\begin{theorem}[Gruppenrekonstruktion aus Einermengenbildern]
\label{b40guide:thm:reconstruction}
Jeder Isomorphismus \(F:\mathcal P_*(G)\to\mathcal P_*(H)\)
bildet die Menge der Einermengen von \(G\) auf diejenige von \(H\) ab.
Es gibt genau einen Gruppenisomorphismus \(\varphi:G\to H\) mit
\[
 F(\{g\})=\{\varphi(g)\}\qquad(g\in G).
\]
Diese Aussage betrifft die Einermengenbilder; sie setzt keine
Erhaltung von Vereinigungen durch \(F\) voraus.
\cite{projectreconstruction,liutringali}
\end{theorem}

\begin{theorem}[Vollständiges Isomorphiekriterium für zwei Gruppen]
Für beliebige Gruppen \(G,H\) gilt
\[
 G\cong H
 \quad\Longleftrightarrow\quad
 (\mathcal P_*(G),\cdot)\cong(\mathcal P_*(H),\cdot).
\]
Jeder Gruppenisomorphismus \(\varphi:G\to H\) induziert dabei den
Isomorphismus \(A\mapsto\varphi[A]\) der Mengenproduktstrukturen.
Bei einem beliebig vorgegebenen Isomorphismus dieser Strukturen
liefert Satz~\ref{b40guide:thm:reconstruction} den zugehörigen
Gruppenisomorphismus über die Einermengen.
\cite{projectreconstruction,liutringali}
\end{theorem}
